% ===== PRÉAMBULE CANONIQUE — à coller VERBATIM en tête de CHAQUE .tex =====
\documentclass[11pt,a4paper]{article}
\usepackage[utf8]{inputenc}\usepackage[T1]{fontenc}\usepackage{lmodern}
\usepackage[french]{babel}
\usepackage{amsmath,amssymb,mathtools}\usepackage{icomma}
\usepackage{siunitx}\sisetup{locale=FR}  % \qty{}{}, \si{}, \ang{}
\usepackage{xcolor}\usepackage{colortbl}
\usepackage{tikz}\usepackage{pgfplots}\pgfplotsset{compat=1.18}
\usepackage[siunitx,europeanresistors]{circuitikz} % circuits (élec.)
\usepackage[most]{tcolorbox}\usepackage{enumitem}\usepackage{array,booktabs}
\usepackage[a4paper,margin=2.2cm]{geometry}
\usetikzlibrary{arrows.meta,calc,angles,quotes,positioning,patterns,%
  backgrounds,shapes.geometric,decorations.pathmorphing,decorations.markings,intersections}
\definecolor{primary}{RGB}{222,49,99}\definecolor{primarydark}{RGB}{150,22,55}
\definecolor{defcol}{RGB}{0,114,178}\definecolor{methcol}{RGB}{52,92,148}
\definecolor{excol}{RGB}{0,135,90}\definecolor{attncol}{RGB}{200,40,55}
\tcbset{boxbase/.style={enhanced,breakable,boxrule=0.9pt,arc=2.5pt,
  left=3mm,right=3mm,top=2.3mm,bottom=2.3mm,fonttitle=\bfseries,coltitle=white}}
% --- boîtes TUTORIEL (noms EXACTS, ne pas renommer) ---
\newtcolorbox{cle}[1][]{boxbase,colback=defcol!6,colframe=defcol,
  title={\textsc{À retenir}\ifx\relax#1\relax\relax\else\textmd{ : \itshape #1}\fi}}
\newtcolorbox{methode}[1][]{boxbase,colback=methcol!7,colframe=methcol,sharp corners,
  title={\textsc{Méthode}\ifx\relax#1\relax\relax\else\textmd{ --- \itshape #1}\fi}}
\newtcolorbox{exemplebox}[1][]{boxbase,colback=excol!5,colframe=excol,
  title={\textsc{Exemple}\ifx\relax#1\relax\relax\else\textmd{ : \itshape #1}\fi}}
\newtcolorbox{piege}[1][]{boxbase,colback=attncol!6,colframe=attncol,
  title={\textsc{Piège}\ifx\relax#1\relax\relax\else\textmd{ : \itshape #1}\fi}}
\newtcolorbox{remarque}[1][]{boxbase,colback=black!4,colframe=black!45,
  title={\textsc{Remarque}\ifx\relax#1\relax\relax\else\textmd{ : \itshape #1}\fi}}
% --- boîtes EXERCICE / CORRIGÉ (noms EXACTS, auto-comptées) ---
\newtcolorbox[auto counter]{exo}[1][]{boxbase,colback=primary!4,colframe=primary,
  title={\textsc{Exercice}~\thetcbcounter\ifx\relax#1\relax\relax\else\textmd{ --- \itshape #1}\fi}}
\newtcolorbox[auto counter]{corrige}[1][]{boxbase,colback=excol!4,colframe=excol,
  title={\textsc{Corrigé}~\thetcbcounter\ifx\relax#1\relax\relax\else\textmd{ --- \itshape #1}\fi}}
\newcommand{\vv}[1]{\vec{#1}}\newcommand{\ds}{\displaystyle}
\newcommand{\R}{\mathbb{R}}\newcommand{\N}{\mathbb{N}}\newcommand{\Z}{\mathbb{Z}}
% (un \usepackage{fancyhdr}+en-tête est possible ; il est ignoré côté web)
% =========================================================================
\begin{document}
% THEME_TITLE: Potentiel, tension et travail électrique
\sisetup{per-mode=symbol}

\begin{center}\large\bfseries\color{primarydark}
Thème 5 — Potentiel, tension et travail électrique
\end{center}

Déplacer une charge dans un champ électrique coûte (ou fournit) de l'énergie. On décrit cela par le
\textbf{potentiel} $V$, la \textbf{différence de potentiel} $U$ et le \textbf{travail} $W$.

\begin{cle}[potentiel et différence de potentiel]
Le \textbf{potentiel} en $M$ est l'énergie par unité de charge pour amener une charge depuis
l'infini : $V_M = W_{\infty\to M}/q$ (scalaire, en volts). Seule la \textbf{différence} se mesure :
\[ \boxed{\,U_{AB} = V_A - V_B\,}\qquad \text{avec}\quad \SI{1}{\volt} = \SI{1}{\joule\per\coulomb} = \SI{1}{\newton\metre\per\coulomb}. \]
\end{cle}

\begin{cle}[travail de la force électrique]
Pour déplacer une charge $q$ (algébrique) de $B$ vers $A$ :
\[ \boxed{\,W_{B\to A} = q\,(V_A - V_B) = q\,U_{AB}\,}\qquad(\si{\joule}). \]
$W>0$ : la force électrique \emph{fournit} du travail (mouvement naturel) ; $W<0$ : il faut fournir
un travail \emph{contre} la force. Le travail est \textbf{proportionnel à la charge} déplacée.
\end{cle}

\begin{cle}[condensateur : champ et potentiel]
Le champ est uniforme et le potentiel décroît \emph{linéairement} de l'armature $+$ vers l'armature
$-$ :
\[ \boxed{\,E = \dfrac{|U_{AB}|}{d}\,}\qquad(\si{\volt\per\metre}). \]
Au milieu des plaques, le potentiel vaut $(V_A+V_B)/2$ : $U_{AP} = U_{AB}/2$.
\end{cle}

\begin{methode}[particule en équilibre entre deux plaques]
Une particule chargée \emph{immobile} entre deux armatures est en équilibre : la force électrique
compense le poids.
\begin{enumerate}[leftmargin=1.6em,topsep=1pt,itemsep=1pt]
  \item \textbf{Signe} : la force $\vv{F}=q\vv{E}$ doit être dirigée \emph{vers le haut} (opposée au
        poids) $\Rightarrow$ compare le sens de $\vv{E}$ à celui de $\vv{F}$ pour trouver le signe de $q$.
  \item \textbf{Norme} : $|q|\,E = m\,g$, avec $E = U/d$, d'où $|q| = \dfrac{m\,g}{E}$.
  \item \textbf{Électrons} : $N = |q|/e$.
\end{enumerate}
\end{methode}

\begin{center}
\begin{tikzpicture}[scale=0.85,>=Stealth,font=\scriptsize]
  \fill[attncol!35] (-2.4,1.5) rectangle (2.4,1.65); \node[attncol,font=\bfseries] at (2.95,1.57){$+$ ($A$)};
  \fill[defcol!35] (-2.4,-1.65) rectangle (2.4,-1.5); \node[defcol,font=\bfseries] at (2.95,-1.57){$-$ ($B$)};
  \foreach \x in {-1.6,-0.8,0,0.8,1.6}{\draw[->,excol,line width=0.8pt] (\x,1.35)--(\x,-1.35);}
  \node[excol] at (-1.9,0.55){$\vv{E}$};
  \node[circle,fill=defcol,minimum size=6mm,inner sep=0pt,text=white,font=\bfseries\scriptsize] (q) at (0,0) {$-$};
  \draw[->,defcol,line width=1.2pt] (0,0.15)--(0,1.15) node[right,font=\scriptsize]{$\vv{F}=q\vv{E}$};
  \draw[->,black,line width=1.2pt] (0,-0.15)--(0,-1.15) node[right,font=\scriptsize]{$\vv{P}=m\vv{g}$};
  \draw[<->,black!55] (-2.55,1.5)--(-2.55,-1.5) node[midway,left]{$d$};
\end{tikzpicture}
\end{center}

\begin{piege}[énergie faible $\ne$ tension faible]
Comme $U = W/q$, une \emph{petite} énergie sous une charge minuscule peut cacher une \emph{grande}
tension. Il faut toujours raisonner sur la d.d.p., pas sur l'énergie brute. À l'inverse, à d.d.p.
fixée, doubler la charge \emph{double} le travail ($W = qU$).
\end{piege}

\end{document}
